\chapter{Implementácia hardvérovej časti}
\label{ImplementaciaHW}

Táto kapitola popisuje hardvérovú realizáciu navrhnutého riadiaceho
systému. Jednotlivé uzly sú rozdelené podľa svojej funkčnej
zodpovednosti — centrálna jednotka zabezpečuje riadiacu logiku
a~multimediálny výstup, distribuované uzly na báze ESP32 zabezpečujú
ovládanie motorov, relé a~senzorickú detekciu vstupu návštevníka.

\section{Centrálna jednotka -- Raspberry Pi}
\label{CentralnaJednotka}

Úlohu hlavného komunikačného a multimediálneho uzla plní
Raspberry Pi (viď kapitola~\ref{sbc-platformy}). Z hľadiska fyzickej
inštalácie bol umiestnený v tesnej blízkosti audiovizuálnej techniky
(priamo za obrazovkou dispeja). Toto riešenie minimalizuje dĺžku
káblových zväzkov, keďže Raspberry Pi disponuje HDMI portom ale aj audio výstup cez 3,5\,mm
jack. Pre integráciu do systému postačuje len bežný Ethernet kábel pre prístup do LAN siete.

Na napájanie Raspberry Pi je využívaný USB-C adaptér - 5,0\,V / 3,0\,A.

\vspace{1cm}
\begin{center}
    \textit{[MIESTO PRE OBRÁZOK: Fotografia zapojeného Raspberry Pi v~rozvádzači
    s~HDMI a~LAN káblovou konektivitou]}
\end{center}
\vspace{1cm}

\section{Distribuované uzly -- ESP32}
\label{DistribuovaneUzly}

Riadiaca logika expozície bola decentralizovaná do niekoľkých
samostatných uzlov postavených na architektúre mikrokontrolérov
rodiny ESP32 (viď kapitola~\ref{mikrokontrolery}). Tento prístup
fyzicky rozdeľuje zodpovednosť do funkčných blokov a~znižuje potrebu
viesť káblové zväzky naprieč miestnosťou. Všetky uzly komunikujú s Raspberry Pi cez Wi-Fi sieť.

Riadiace a~silové uzly (pre ovládanie motorov a~spínanie relé)
boli umiestnené spoločne vo vnútri inštalačného rozvádzača
integrovaného priamo do fyzického exponátu expozície. Senzorický
tlačidlový uzol pre používateľský vstup bol namontovaný oddelene
pri vstupe do miestnosti.

\subsection{Motorický uzol -- riadenie DC motorov}
\label{UzolMotorov}

Pre ovládanie pohonov exponátov bola zvolená dvojica komponentov
ESP32 DevKit V1 a~výkonový budič BTS7960, ktorých kombinácia
pokrýva celý rozsah požiadaviek od riadiacej logiky až po silové
prúdové špičky pri rozbehoch. Konkrétne fyzické exponáty, ktoré
tento uzol poháňa, sú popísané v~kapitole~\ref{MechanickaKonstrukcia}.

Hardvérová implementácia uzla pozostáva z~nasledujúcich komponentov:
\begin{itemize}
    \item \textbf{Mikrokontrolér ESP32 DevKit V1:} Funguje ako mozog uzla,
        ktorý spracováva príkazy z MQTT brokera a tým generuje PWM signál
        pre ovládanie motorov.
    \item \textbf{2x Výkonový budič BTS7960:} Polovodičový H-mostík slúži na
        riadenie výkonu motorov (viď kapitola~\ref{dc-motory}). Tento masívne
        dimenzovaný obvod bezpečne odoláva vysokým prúdovým špičkám pri
        rozbehoch (až do 43\,A).
\end{itemize}

\vspace{1cm}
\begin{center}
    \textit{[MIESTO PRE OBRÁZOK: Schéma zapojenia uzla pre DC Motory
    (ESP32, BTS7960, 12V zdroj)]}
\end{center}
\vspace{1cm}

\subsection{Reléový uzol -- spínanie záťaží}
\label{UzolRele}

Na spínanie záťaží ostatných prvkov expozície slúži vyhradený reléový uzol:
\begin{itemize}
    \item \textbf{Priemyselný I/O modul Waveshare:} Základom tohto uzla je
        komerčný hotový modul s~čipom \textbf{ESP32-S3}. Obsahuje integrovanú
        podporu pre \textit{Power over Ethernet} (PoE) a~zbernicu RS485.
    \item \textbf{10\,A mechanické relé:} Modul disponuje ôsmimi nezávislými
        relé. Spínajú priamo striedavých
        230\,V na chod dekoratívnej \textbf{Edisonovej žiarovky},
        \textbf{výrobníka dymu} (ktorý vháňa paru do rozvedeného potrubia)
        a~tiež prúd pre \textbf{svetlo simulujúce oheň} v~kotle. Ďalšou
        vyhradenou sekciou je spínanie jednosmernej 12\,V vetvy určenej na
        napájanie podsvietenia dvoch oddelených skupín \textbf{industriálnych
        analógových budíkov} a~skrytých dizajnových LED pásov.
    \item \textbf{Dôvod nasadenia modulu:} Namiesto lokálne osadených
     polovodičových relé (SSR) (viď kapitola~\ref{spinanie-sietoveho-napatia})
        bola na zaistenie požiarnej a~elektrickej bezpečnosti použitá táto
        certifikovaná hotová doska. V~prípade hardvérového zlyhania v múzejnej
        prevádzke umožňuje servisným technikom okamžitú, jednoduchú
        a~bezpečnú výmenu \uv{kus za kus}.
\end{itemize}

\vspace{1cm}
\begin{center}
    \textit{[MIESTO PRE OBRÁZOK: Schéma zapojenia relé uzla
    (Waveshare ESP32-S3, 230V obvody a 12V obvody)]}
\end{center}
\vspace{1cm}

\subsection{Vstupný tlačidlový uzol}
\label{UzolSenzorov}

Fyzický vstup zameraný na interakciu s~návštevníkmi je spracovaný
samostatným uzlom:
\begin{itemize}
    \item \textbf{Samostatný mikrokontrolér:} Na vstupnom mieste pred
        stenou je fyzicky osadená inštalačná krabička. Aby si
        zachovala kompaktné vonkajšie zapojenie, disponuje vlastným
        DC napájacím konektorom pre privedenie 5\,V napätia k~vloženej
        doske \textbf{ESP32 DevKit V1}.
    \item \textbf{Antivandal tlačidlo:} Kvôli zabezpečeniu odolnosti
        pred opakovaným a~tvrdým návštevníckym stláčaním bol zvolený
        priemyselný \textbf{19\,mm vodotesný kovový spínač}, ktorý je
        prepojený s~krabičkou riadiaceho modulu pomocou leteckého
        konektora typu \textbf{GX12 (4-pin)}. Cez tento štvoritý
        konektor sú súčasne vedené všetky vodiče uzla: signálový
        vodič tlačidla, napájanie LED prstenca, uzemnenie
        a~referenčné 5\,V. Spínač slúži ako prvotný štartovací
        impulz pre chod interaktívnej scény.
    \item \textbf{Riadenie LED prstenca cez MOSFET spínač:} Záťaž
        LED prstenca bola namiesto priameho napájania z~GPIO pinu
        riadená cez kompaktný MOSFET PWM spínací modul
        s~tranzistorom F5305S (viď kapitola~\ref{dc-motory}), čím
        bola výstupná periféria mikrokontroléra ochránená pred
        prúdovým preťažením a~zároveň bolo umožnené PWM stmievanie
        prstenca.
\end{itemize}

\vspace{1cm}
\begin{center}
    \textit{[MIESTO PRE OBRÁZOK: Fotografia inštalačnej
    krabičky tlačidlového uzla s~viditeľným MOSFET modulom, ESP32
    doskou a~GX12 konektorom]}
\end{center}
\vspace{1cm}

\vspace{1cm}
\begin{center}
    \textit{[MIESTO PRE OBRÁZOK: Schéma zapojenia vstupného
    tlačidlového uzla -- ESP32, externý pull-up rezistor na
    GPIO\,32, MOSFET F5305S pre napájanie LED prstenca
    a~GX12 konektor]}
\end{center}
\vspace{1cm}

\section{Mechanická konštrukcia exponátov}
\label{MechanickaKonstrukcia}

Fyzické exponáty inštalácie CHRONOS boli navrhnuté
v~prostredí \textbf{Fusion 360}. Celková sústava
obsahuje \textbf{24 ozubených kolies} z~preglejky
s~jednotným modulom $m = 8$. Pre každé koleso bol
vytvorený samostatný 3D model; pri menších kolesách
bol G-kód generovaný a~obrobky boli vyfrézované
na dostupnej CNC fréze. Väčšie kolesá boli zadané
externej firme, ktorá disponuje strojom s~väčším
pracovným priestorom.

Súčasťou návrhu bol aj \textbf{dekoratívny piest}
pripevnený na sústave kolies, ktorý vykonáva
pohyb bez mechanického zaťaženia, a~\textbf{hodinový
mechanizmus} s~priemerom 1{,}3\,m vrátane
planetového prevodu s~vypočítaným počtom zubov
pre požadovaný prevodový pomer. Montáž
a~povrchovú úpravu exponátov zabezpečila
externá firma. Podrobná konštrukčná dokumentácia
(výkresy, modely, prevodové výpočty) presahuje
rámec tejto práce a~nie je jej predmetom.